\documentclass[11pt]{article}

\usepackage[margin=1in]{geometry}
\usepackage{amsmath,amssymb}
\usepackage{hyperref}
\usepackage{listings}
\usepackage{xcolor}
\usepackage{booktabs}
\usepackage{enumitem}

\lstset{
  basicstyle=\ttfamily\small,
  columns=fullflexible,
  keywordstyle=\bfseries,
  breaklines=true,
  frame=single,
  language=C++,
}

\newcommand{\B}{\mathcal{B}}

\title{Continuous Cuts on the Product Bernoulli Manifold}
\author{Riemannian Multilevel Optimization (RMO) framework}
\date{}

\begin{document}
\maketitle

\section{Overview}

This note describes the \texttt{rmo::cc} module, which implements the binary
continuous-cuts image segmentation problem from Section 6.3 of the paper inside the
existing \texttt{rmo} framework. It follows the same split already used for the
Gross--Pitaevskii problem (\texttt{rmo::gpe}): the problem-independent core in
\texttt{rmo::ropt} (manifolds, transfer operators, oracles, the multilevel scheme) is
reused unchanged, and \texttt{rmo::cc} supplies the parts specific to continuous cuts --
the geometry of the product Bernoulli manifold (Table 9), the pixel-grid discretization
and objective (\S 6.3.1), the grid transfer operators (\S 6.3.2, Table 10), and the
two-level coarse model.

Every module in Table~\ref{tab:status} is implemented and tested on the $41\times41$
synthetic disk this module's drivers use, including a two-level driver wiring the coarse
model into the framework's generic multilevel solver (\texttt{FullApproximationScheme}); an
earlier revision of this note reported that the driver did not demonstrate a clear speedup
over the single-level baseline there. That conclusion was wrong -- it was caused by two
parameter choices in \texttt{main\_cc\_coarse.cc} (an over-damped coarse line search and a
missing lockout on consecutive coarse corrections), not by anything structural, and is
retracted in \S\ref{sec:mlcc-design}; with those two parameters corrected the driver reaches
in one fine iteration an energy the single-level baseline needs 16 iterations to reach,
matching the Python reference on the same problem. The port can now also run the paper's
own $960\times1280\to240\times320$ cow-image problem (\S6.3): \texttt{BernoulliGridTransfer}
was generalized to even fine sizes and a new \texttt{ComposedGridTransfer}/
\texttt{ComposedVectorTransport} pair composes the two factor-2 average-pool steps the
paper's coarse level spans, with the actual JPEG decode/resize/blur delegated to a
preprocessing script (\texttt{test/cc/prepare\_cow\_data.py}) that reuses the reference's own
skimage-based pipeline rather than reimplementing it. \texttt{main\_cc\_cow.cc} matches the
reference to 6 digits at the checked iterations and shows a genuine $3.16\times$ CPU-time
speedup to 90\% convergence -- the first problem size in this document where the two-level
driver is actually faster -- alongside a $3\times$ CPU-time \emph{loss} by 300 iterations,
caused by the coarse-correction trigger firing on 145 of 150 eligible iterations at this
scale (see \texttt{REVIEW-continuous-cuts.md} \S8 for the full account, including why this is
the same ratio the retracted paragraph above blamed for the synthetic disk's stall, now
real rather than an artifact). Options 2/3/5 and a 3+ level hierarchy remain out of scope.

Every non-trivial formula here has been checked two ways: once as an internal consistency
check (e.g.\ a gradient against a central difference of the matching value), and once
against the paper's own reference implementation, \texttt{RMO-continuous-cuts}
(\S\ref{sec:verification}). This caught several real bugs along the way -- a wrong
coarse-gradient formula, a wrong description of which Table 10 options coincide, the
retraction reaching the exact domain boundary in floating point, and two oracle instances
silently losing track of each other's updates -- each noted where it was found, since how
it was caught matters more than the bug itself.

\begin{table}[h]
\centering
\begin{tabular}{@{}lll@{}}
\toprule
Component & File & Status \\
\midrule
Pixel grid + forward-difference operator ($D$)   & \texttt{cc/grid.h}            & done \\
Fisher--Rao metric                               & \texttt{cc/metric.h}          & done \\
Manifold geometry (Table 9)                      & \texttt{cc/manifold.h}        & done \\
Objective (ambient value/gradient)               & \texttt{cc/cc.h}              & done \\
Fine-level oracle                                & \texttt{cc/oracle.h}          & done \\
Grid transfer (injection, bilinear)              & \texttt{cc/interpolate.h}     & done \\
Point map + vector transport (Table 10)          & \texttt{cc/transport.h}       & done \\
Coarse-model oracle                              & \texttt{cc/oracle\_coarse.h}  & done \\
Single-level driver                              & \texttt{src/main\_cc.cc}      & done \\
Two-level driver (\S\ref{sec:mlcc-design})       & \texttt{src/main\_cc\_coarse.cc} & done \\
Tests (4 files, \S\ref{sec:testing})             & \texttt{test/cc/*.cc}         & done \\
\bottomrule
\end{tabular}
\caption{Status of the \texttt{rmo::cc} module.}
\label{tab:status}
\end{table}

\section{Design}

\subsection{Simplifications relative to Gross--Pitaevskii}

Two things make this module noticeably simpler than \texttt{rmo::gpe}:

\begin{itemize}[nosep]
  \item \textbf{No linear solves.} The Fisher--Rao metric is diagonal,
        $G(\phi) = \operatorname{diag}(1/(\phi(1-\phi)))$, and every formula in Table 9 is a
        closed-form elementwise expression. There is no mass matrix, and nothing like GP's
        \texttt{InverseMatrix}/\texttt{PreconditionInverse} is needed anywhere. The only
        matrix-shaped objects in the module are the sparse gradient operator $D$ and the
        grid transfer operators.
  \item \textbf{Plain array indexing instead of finite elements.} The grid transfer
        operators -- injection $J_h^H$, bilinear interpolation $B_H^h$, and full-weighting
        $F_h^H=\tfrac14(B_H^h)^\top$ -- do not need a finite-element mesh at all. The
        reference implementation's \texttt{primitives.py} builds them directly as
        index/parity arithmetic on the pixel array (e.g.\ \texttt{phi[::2,::2]} for
        injection), and \texttt{cc/interpolate.h} does the same: given a coarse and a fine
        pixel grid related by one exact doubling per dimension, each operator is a small
        case split on whether a fine index is even or odd in each coordinate. An earlier
        version of this module instead reused \texttt{rmo::fe::LinearTransferMG}
        (deal.II's finite-element mesh transfer, built for GP), which works but pulls in a
        whole mesh library for what is just array indexing.
\end{itemize}

\subsection{Files}

\paragraph{\texttt{grid.h}: pixel grid and the $D$ operator.} \texttt{PixelGrid(rows, cols)}
represents an image as an $\mathrm{FE\_Q}(1)$ finite-element space with one degree of
freedom per pixel (chosen so later grid-transfer code, \S\ref{sec:interpolate-design}, has a
DoFHandler to work with if needed, not because any finite-element machinery is used for the
operators themselves). Since deal.II does not number degrees of freedom in raster (row-major)
order, \texttt{PixelGrid} builds the permutation between raster and DoF order once, from the
degrees of freedom's physical locations. All vectors ($\phi$, $\rho$, gradients) live in DoF
order everywhere in the module; \texttt{to\_dof\_order}/\texttt{to\_raster\_order} are the
only place raster order appears, at the boundary with image data.

\texttt{ForwardDifference} builds the sparse $2n\times n$ matrix $D=[D_1;D_2]$ of eq.~(42):
one-sided forward differences, with a zero row for pixels that have no forward neighbor
(the last column for $D_1$, the last row for $D_2$). Row $i$ of $D_1$, row $i$ of $D_2$, and
column $i$ of the identity all refer to the same pixel $i$ in the same DoF order used for
$\phi$ everywhere else.

\paragraph{\texttt{metric.h} and \texttt{manifold.h}: Table 9.} The Fisher--Rao metric
(inner product, norm, and the diagonal weight $\phi(1-\phi)$ they are built from) lives in
its own namespace, \texttt{metric::fisher\_rao}, separate from the retraction and lifting
map in \texttt{manifold.h}. This keeps ``the metric'' and ``the manifold'' as two
independent, separately reusable things -- for instance \texttt{oracle.h}'s
\texttt{norm}/\texttt{inner} only need the metric, not the retraction.

The retraction $R_\phi$ and lifting map $\mathcal L_\phi$ (Table 9) both reduce to a
logit-space identity:
\[
  R_\phi(v) = \sigma\big(\operatorname{logit}(\phi) + G(\phi) v\big), \qquad
  \mathcal{L}_\phi(z) = G(\phi)^{-1}\big(\operatorname{logit}(z) - \operatorname{logit}(\phi)\big),
\]
where $\sigma$ is the logistic sigmoid, $\operatorname{logit}$ its inverse, and
$G(\phi)v=v/(\phi(1-\phi))$. Implementing them this way (through stable
\texttt{sigmoid}/\texttt{logit} functions) instead of the raw exponential form written in
the table avoids overflow as $\phi$ approaches 0 or 1, and is exactly the identity the
paper itself uses (p.~29) to simplify the coarse-model correction term (\S\ref{sec:oracle-coarse-design}).

Differentiating gives $\mathrm D R_\phi(v)[w] = z(1-z)\,w/(\phi(1-\phi))$ (with
$z=R_\phi(v)$) and $\mathrm D\mathcal L_\phi(z)[u] = (\phi(1-\phi)/(z(1-z)))\,u$. The latter
is a diagonal linear map, so it is its own adjoint under the plain Euclidean dot product --
but \emph{not} under the Fisher--Rao metric, which matters for the coarse gradient
(\S\ref{sec:oracle-coarse-design}) and is easy to get wrong, since GP's ellipsoid manifold
uses a metric that \emph{is} its own metric-adjoint under an argument swap, and this one is
not.

\paragraph{\texttt{cc.h}: the objective.} \texttt{ContinuousCutsFunctional} stores $\rho$
(the data term; building it from an image and foreground/background seed patches, as in
Fig.~12, is left to the driver), the $D$ operator, $\alpha$, $\varepsilon$, and implements:
\begin{itemize}[nosep]
  \item \texttt{value(phi)}: $\rho^\top\phi + \alpha\sum_i\sqrt{(D_1\phi)_i^2+(D_2\phi)_i^2+\varepsilon^2}$;
  \item \texttt{gradient(phi, dst)}: the \emph{ambient} (Euclidean) gradient, not yet the
        Riemannian one -- see the derivation below;
  \item \texttt{directional\_derivative(phi, z)}, reusing the same evaluation.
\end{itemize}
\texttt{update(phi)} caches $D\phi$ so the three methods above do not each reapply $D$; like
GP's \texttt{GrossPitaevskiiFunctional}, they assume \texttt{phi} matches the point passed
to the most recent \texttt{update}.

\paragraph{Ambient gradient, and converting it to a Riemannian one.} These derivations are
short enough to write out. Let $g_i(\phi)=\sqrt{(D_1\phi)_i^2+(D_2\phi)_i^2+\varepsilon^2}$
be the $i$-th smoothed-TV term, so $E_\varepsilon^{CC}(\phi)=\rho^\top\phi+\alpha\sum_i g_i(\phi)$.
Since $D_1,D_2$ are linear,
\[
  \frac{\partial g_i}{\partial\phi_k} = \omega_i\big((D_1)_{ik}(D_1\phi)_i+(D_2)_{ik}(D_2\phi)_i\big),
  \qquad \omega_i=\frac1{g_i(\phi)},
\]
and summing over $i$ is a $D^\top$-contraction with the weighted vector, giving the ambient
gradient
\[
  g := \operatorname{grad}_{\mathrm{amb}}E_\varepsilon^{CC}(\phi) = \rho + \alpha D^\top(I_2\otimes\operatorname{diag}(\omega))D\phi.
\]
This is exactly what \texttt{cc.h::gradient} computes. Since $\B=(0,1)^n$ is an open subset
of $\mathbb R^n$, its tangent space is just $\mathbb R^n$, and the directional derivative is
the plain dot product $\mathrm DE_\varepsilon^{CC}(\phi)[v]=g^\top v$. The Riemannian
gradient is defined by asking for the same identity under the Fisher--Rao pairing instead,
$\langle\operatorname{grad}E_\varepsilon^{CC}(\phi),v\rangle_\phi = g^\top v$ for every $v$;
writing out $\langle u,v\rangle_\phi=u^\top G(\phi)v$ turns this into one linear equation,
solved by
\[
  \operatorname{grad}E_\varepsilon^{CC}(\phi) = G(\phi)^{-1}g = \phi(1-\phi)\odot g,
\]
which is Table 9's gradient row. Because $G(\phi)$ is diagonal this ``solve'' is just an
elementwise multiply (\texttt{detail::grad\_fisher\_rao} in \texttt{oracle.h}) -- unlike the
mass- or energy-metric conversions in \texttt{rmo::gpe}, which need an actual linear solve.
The same equation, read in the other direction, is why \texttt{oracle.h}'s
\texttt{apply\_metric} is $G(\phi)$ itself (divide by $\phi(1-\phi)$, not multiply): it maps
a tangent vector to the covector it defines, i.e.\ ambient-to-Riemannian run backwards.

\paragraph{\texttt{oracle.h}: the fine-level oracle.} \texttt{FisherRaoOracle} computes the
Riemannian gradient as above, with no linear solve. \texttt{norm}/\texttt{inner} call
\texttt{metric::fisher\_rao} directly, and since those (per the \texttt{OracleBase}
interface) take no base point of their own, \texttt{update(phi)} caches \texttt{phi} for
them to use -- the same role a (here state-independent) mass-matrix reference plays for
GP's \texttt{MassOracle}. The paper uses the Fisher--Rao metric throughout for this problem
(unlike GP's Mass/Energy/Frobenius choice), so this is the only oracle class needed, plus a
small \texttt{ContinuousCutsResidual} helper for the stationarity measure
$\|\operatorname{grad}E_\varepsilon^{CC}(\phi)\|_\phi$ used as the residual in Figs.~13--14.
A \texttt{FISHER\_RAO} entry was added to the shared \texttt{MetricKind} enum
(\texttt{option\_types.h}) for this. It is deliberately \emph{not} added to
\texttt{option.h}'s CLI-parseable list: that list feeds \texttt{main\_coarse.cc}'s
\texttt{--metric} flag, which is GP-specific and has no handler for it, so parsing it from
the command line would only reach a crash further down. Add it there once a CC driver with
its own \texttt{--metric} handling exists.

\paragraph{\texttt{interpolate.h}: grid transfer.}
\label{sec:interpolate-design}
\texttt{BernoulliGridTransfer} holds a coarse and a fine \texttt{PixelGrid} related by one
exact doubling per dimension ($\mathrm{fine.rows()}=2\cdot\mathrm{coarse.rows()}-1$, and
likewise for columns -- the same relationship one step of mesh refinement gives). Every
fine index falls into exactly one of three cases relative to its coarse neighbors: both
coordinates even (it coincides with a coarse pixel), one odd (it sits at a coarse edge
midpoint), or both odd (a coarse cell center). These are the same three cases
\texttt{test\_operator.cc} documents for bilinear interpolation on nested meshes, and under
the exact-doubling assumption the neighboring coarse index is always in range, so no
special-casing is needed at the image boundary.
\begin{itemize}[nosep]
  \item \texttt{to\_coarse\_mesh} is injection, $J_h^H$: $\mathrm{coarse}(I,J)=\mathrm{fine}(2I,2J)$.
  \item \texttt{to\_fine\_mesh} is bilinear interpolation, $B_H^h$, using the three-case
        weights above (1; $\tfrac12,\tfrac12$; $\tfrac14,\tfrac14,\tfrac14,\tfrac14$).
  \item \texttt{Tfine} is the exact transpose of \texttt{to\_fine\_mesh}: the same three
        cases, written as a scatter (accumulate into the coarse array) instead of a gather.
        Since $F_h^H=\tfrac14(B_H^h)^\top$, \texttt{Tfine} \emph{is} $4F_h^H$, so
        full-weighting is simply \texttt{Tfine} scaled by $\tfrac14$ wherever it is needed
        (\S\ref{sec:transport-design}) -- no separate stencil, the same
        transpose-as-adjoint pattern \texttt{lac.h}'s \texttt{TransposeMatrix} already uses
        elsewhere in the codebase.
\end{itemize}
No metric weighting appears in this file; it is pure grid arithmetic. The Fisher--Rao
weighting eq.~(43) asks for is applied by the code that calls these operators
(\S\ref{sec:transport-design}), not by the operators themselves.

\paragraph{\texttt{transport.h}: point map and vector transport (Table 10).}
\label{sec:transport-design}
\texttt{BernoulliPointTransfer} implements the point restriction/prolongation $r,p$ from
\S 6.3.2. $p$ always uses bilinear interpolation in logit coordinates,
$p(\psi)=\sigma(B_H^h\operatorname{logit}(\psi))$. $r$ is chosen at construction between
plain injection (calling \texttt{to\_coarse\_mesh} directly, with no logit round trip --
injection commutes with any elementwise map, checked in \texttt{test\_interpolate.cc}) and
full-weighting, $r(\phi)=\sigma(\tfrac14\texttt{Tfine}(\operatorname{logit}(\phi)))$.

Table 10 lists five vector-transport options; comparing them against the reference's
\texttt{operators.py} shows Options 1 and 4 use \emph{identical} $P$ and $R$ formulas and
differ only in which restriction builds $\psi=r(\phi)$ (injection vs.\ full-weighting) --
one class, \texttt{GeometricBilinearTransport}, covers both:
\[
  P^\phi_\psi(v) = G(\phi)^{-1}\big(B_H^h\,G(\psi)v\big), \qquad
  R^\psi_\phi(v) = (B_H^h)^\top v = 4F_h^H v \ \text{(via \texttt{Tfine}, unweighted)}.
\]
Choosing Option 1 vs.\ Option 4 is then just a choice of
\texttt{BernoulliPointTransfer::Restriction}, not a different transport class. Options 2, 3,
5 are not implemented: 2 and 5 skip the $P$-side weighting and underperform in the paper's
own results (Figs.~13--14), and 3 needs a Moore--Penrose pseudoinverse (an
$(FF^\top)^{-1}$ solve) for an option the paper also finds triggers coarse corrections far
too often -- not worth the extra complexity given 1 and 4 already win.

\paragraph{\texttt{oracle\_coarse.h}: the coarse model.}
\label{sec:oracle-coarse-design}
\texttt{CoarseOracleBase} (already part of \texttt{rmo::ropt}) needs no changes; one new
class, \texttt{ContinuousCutsCoarseOracle}, plays both roles
\texttt{FullApproximationScheme} expects from a coarse model (again because only one metric
is used, unlike GP's several). Because
$G_h(\phi)=\mathrm D\operatorname{logit}(\phi)$, the correction term in the coarse model
$q_k(z) = E_H(z) - \langle w_k,\mathcal L_{\psi_k}(z)\rangle_{\psi_k}$ collapses to the
affine, metric-free expression $w_k^\top(\operatorname{logit}(z)-\operatorname{logit}(\psi_k))$
(paper, p.~29), so
\[
  q_k(z) = E_H(z) - w_k^\top\big(\operatorname{logit}(z)-\operatorname{logit}(\psi_k)\big), \qquad
  \mathrm Dq_k(z)[v] = \mathrm DE_H(z)[v] - w_k^\top\Big(\frac{v}{z(1-z)}\Big).
\]

Its Riemannian gradient took a second attempt to get right, and the mistake is worth
recording. Writing $\langle\operatorname{grad}q_k(z),v\rangle_z=\mathrm Dq_k(z)[v]$ and
$g=\operatorname{grad}_{\mathrm{amb}}E_H(z)$, the right-hand side is
$g^\top v - w_k^\top G(z)v = (g-G(z)w_k)^\top v$, so
\[
  \operatorname{grad}q_k(z) = \operatorname{diag}(z(1-z))\,g \;-\; w_k.
\]
The first version of this code instead computed $\operatorname{diag}(z(1-z))(g-w_k)$,
i.e.\ it also rescaled $w_k$, which is wrong. The paper's general formula,
$\operatorname{grad}q(z)=\operatorname{grad}f_H(z)-\mathrm D\mathcal L_{\psi_k}(z)^*[w_k]$
(Prop.~3.1), explains why: the adjoint $\mathrm D\mathcal L_{\psi_k}(z)^*$ is taken between
the metrics at $\psi_k$ and at $z$, and for the Fisher--Rao metric this adjoint turns out to
be the \emph{identity} map (checked both by direct calculation and numerically against the
reference, to about $10^{-10}$) -- \emph{not} the Euclidean self-adjoint that
\texttt{bernoulli::retract\_inv\_diff\_adjoint} computes, which answers a different
question (adjoint under the plain dot product, not under a metric that changes from point
to point). The bug was caught by testing against Proposition 3.2 (below), not by re-reading
the derivation -- worth noting since it means the check, not the formula, deserves the
credit.

\texttt{ContinuousCutsCoarseOracle::gradient} gets $g$ by calling \texttt{T\_coarse.gradient}
(the Riemannian gradient) and then \texttt{T\_coarse.apply\_metric} (which divides by
$z(1-z)$, undoing the weighting) -- both already on the \texttt{OracleBase} interface, so no
downcast to the concrete oracle type is needed, unlike GP's coarse oracles.
\texttt{FisherRaoOracle}'s constructor gained an unused \texttt{SolverOptions} parameter
(matching GP's \texttt{FrobeniusOracle}) purely so it satisfies the interface
\texttt{FullApproximationScheme} expects; see \S\ref{sec:solveroptions} for more on this.

\paragraph{\texttt{main\_cc.cc}: single-level driver.} A synthetic image (a bright disk on
a dark background) with a foreground/background seed patch known to lie inside/outside it
stands in for a loaded image and its seed regions (cf.\ Fig.~12); $\rho$ is built from it
exactly as in eq.~(41)/(42), with $\alpha=0.1$ matching the paper's fine-level regularization
weight (\S 6.3.4) and $\varepsilon=10^{-2}=\sqrt{\varepsilon_{\text{ref}}}$
(\S\ref{sec:mlcc-design}'s $\varepsilon$-convention paragraph). It wires
\texttt{FisherRaoOracle} and \texttt{BernoulliManifold} directly into the already
problem-independent \texttt{GradientDescent} solver and \texttt{ConvergenceTableObserver} --
no CC-specific solver code was needed. On a $41\times41$ grid it runs all 300 configured
iterations (the residual plateaus at $1.446\times10^{-5}$ from iteration $\sim\!161$ on and
never reaches \texttt{tol\_residual}$=10^{-6}$; an earlier revision of this note claimed
convergence at 147 iterations, which \texttt{REVIEW-continuous-cuts.md} \S3.5 found does not
match the committed driver's actual output), with energy monotonically from $184.4$ to
$-110.6$ and a reduction rate that drifts rather than settling at a fixed ratio -- i.e.\ visibly
sublinear, matching the paper's own observation that RGD is sublinear for this problem
(\S 6.4).

\paragraph{\texttt{main\_cc\_coarse.cc}: two-level driver.}
\label{sec:mlcc-design}
Same synthetic problem, at a $21\times21$ coarse / $41\times41$ fine resolution (one exact
doubling apart, as \texttt{interpolate.h} requires), wired through
\texttt{FullApproximationScheme} with \texttt{ContinuousCutsCoarseOracle} playing the
coarse-model role and Option 1 (injection restriction, geometric bilinear transport) as the
vector transport. Getting this to run at all surfaced three real bugs, none of them in the
formulas themselves (already checked exhaustively, \S\ref{sec:oracle-coarse-design} and
\S\ref{sec:testing}) but in assumptions the framework makes that \texttt{rmo::cc}'s single-metric,
no-linear-solve design happened not to satisfy:
\begin{itemize}[nosep]
  \item \textbf{Retraction reaching the exact boundary.} \texttt{FullApproximationScheme}
        constructs \texttt{O\_level} and \texttt{T\_level} as two separate oracle instances
        over the same objective, and an early, large step in a line search retracted a
        point so close to $0$ or $1$ that \texttt{sigmoid} rounded to exactly that value in
        double precision -- which then failed \texttt{logit} on the next call, even though
        $\B=(0,1)^n$ is mathematically open and the e-retraction never actually reaches its
        boundary. \texttt{sigmoid} now clamps its output away from the exact endpoints
        (\texttt{metric::fisher\_rao::MIN\_WEIGHT}, matching the reference's own
        \texttt{apply\_G} clamp for the same reason).
  \item \textbf{Two oracle instances losing track of each other.} \texttt{FullApproximationScheme}
        assumes -- its own comment says so -- that calling \texttt{update} on one of two
        oracle instances sharing an objective makes that update visible to the other. This
        holds for GP's oracles because their state-dependent quantities live on the shared
        functional/system object; it did not hold for \texttt{FisherRaoOracle}, which
        cached its own base point privately, so a second instance's \texttt{norm}/\texttt{inner}
        would see whatever it was last individually updated to -- usually nothing.
        \texttt{ContinuousCutsFunctional} now caches the base point itself
        (\texttt{get\_phi()}), and \texttt{FisherRaoOracle} reads it from there -- the
        smallest fix that matches GP's existing convention (state-dependent quantities live
        on the shared functional, oracles are thin views over it), rather than either of two
        heavier alternatives considered: making \texttt{norm}/\texttt{inner}/\texttt{apply\_metric}
        take an explicit base point (removes the hazard entirely, but touches every oracle
        and call site in the framework, including GP's, which do not need it), or having
        \texttt{FullApproximationScheme} share one oracle instance instead of two when they
        are the same type (does not generalize -- GP genuinely needs \texttt{O\_level} and
        \texttt{T\_level} to differ in metric sometimes). \texttt{OracleBase}'s declaration
        (\texttt{ropt/oracle\_base.h}) now states this as an explicit requirement on
        implementers, since the framework already knew about the hazard -- a \texttt{TODO}
        in \texttt{oracle\_coarse\_base.h} says as much -- but had never written it down
        where a new oracle would be written.
  \item \textbf{The coarse model can run away to the boundary.} $q_k(z)$'s correction term
        is $\pm w_k^\top\operatorname{logit}(z)$, unbounded as $z_i\to0$ or $1$; a full
        ($\alpha=1$) first Armijo step can retract right up to the boundary in one shot
        while still satisfying the sufficient-decrease condition, since the objective
        genuinely has no lower bound in that direction. This is a property of the
        construction itself (Prop.~3.1's correction term is linear in $\operatorname{logit}$),
        not a bug to fix in \texttt{oracle\_coarse.h}. An earlier revision of this driver
        damped the coarse solver's initial step size (\texttt{ls.alpha}$=0.01$) to give
        backtracking a chance to engage before a single step could do this; \S\ref{sec:mlcc-design}
        below explains why that damping, rather than the run-away it guarded against, was
        itself the cause of the driver's reported failure to speed up, and why the reference's
        own approach -- \texttt{ls.alpha}$=1$ with boundary clamping instead of step damping
        (\texttt{optimizer.py:15,49}, matching this port's own \texttt{MIN\_WEIGHT} clamp in
        \texttt{metric.h}) -- is used instead.
\end{itemize}

With those fixed, the driver runs to completion without error: energy decreases
monotonically, and with the coarse correction disabled (via an unreachably high $\eta$) it
reproduces \texttt{main\_cc.cc}'s single-level result almost to machine precision, as it
should.

\paragraph{Retraction: the earlier ``no speedup'' conclusion was wrong, and about the wrong
component.} An earlier revision of this section argued at length -- from the ratio
$\|R\operatorname{grad}f_h(x)\|/\|\operatorname{grad}f_h(x)\|$ growing from about 2 early in
the run to over 200 later -- that Option 1's vector transport ($R=4F_h^H$, a summing rather
than averaging restriction) was structurally unable to produce a speedup on this problem: its
DC gain of 4 was said to amplify the gradient's norm once the segmentation's TV term
concentrates at a sharpening interface, decoupling the trigger condition's two sides. That
argument is retracted. \texttt{REVIEW-continuous-cuts.md} \S3.2 and \S4.2 (independent audit,
measurements reproduced by rerunning the committed and modified drivers) trace the ratio's
growth instead to the fine-level trajectory \emph{stalling}, driven by two parameter choices
in this file, neither of them the vector transport:
\begin{itemize}[nosep]
  \item \texttt{options\_coarse.ls.alpha = 0.01} left the coarse model unminimised (10 steps of
        exactly $0.01$, backtracking only -- Armijo never got the chance to try a larger step),
        so the prolonged correction was too small to move $x$ meaningfully.
  \item \texttt{options\_fas.coarse\_every = 1} meant every fine iteration after the trigger
        first fired was itself a candidate for another coarse correction, with no intervening
        fine-level gradient step (violating \S6.4's rule, cf.\ eq.~(17)) -- so a rejected coarse
        step was followed by another rejected coarse step, forever.
\end{itemize}
With \texttt{ls.alpha}$=1$ (matching \texttt{optimizer.py:8}, undamped on every level) and
\texttt{coarse\_every}$=2$ (matching \texttt{multilevel.py:148}'s lockout), the driver reaches
$E=-109.62$ after the \emph{first} fine iteration -- \texttt{main\_cc.cc}'s single-level
baseline needs 16 iterations to reach that energy -- and $E=-112.23$ by iteration 300,
matching single-level to five digits, with only a handful of rejected line searches rather
than the 2464 the unfixed driver produced. The Python reference, run on the identical
synthetic image by \texttt{test/cc/compare\_reference.py} (added alongside this audit), shows
the same one-iteration jump ($E=-108.55$ under this port's $\varepsilon$ convention). The
ratio growing to $\sim\!200$ in the original run was measuring a point the algorithm should
never have stayed at -- the stall's frozen gradient -- not a per-step property of $R$; in the
fixed run the ratio is $\sim\!10^5$ by iteration 3 and the trigger fires every time it is
checked, exactly as in the reference (where only the $\mu$ gate, not the trigger itself, stops
corrections after the first).

A second, unrelated deviation from the reference was tested and ruled out as a contributing
cause, independently of the above. \texttt{main\_cc\_coarse.cc} originally built its coarse
image by injecting the fine one and recomputing $\rho$ at coarse resolution;
\texttt{problem.py} instead computes $\rho$ once at the finest resolution and average-pools
\emph{that} down to each coarser level (the seed patches are applied exactly once). The driver
now matches the reference -- $\rho$ is coarsened directly, via full-weighting
($F_h^H=\tfrac14\texttt{Tfine}$, this grid's analogue of average-pooling) -- and
\texttt{REVIEW-continuous-cuts.md} \S3.3 confirms the effect is negligible on this image
($-106.94$ vs.\ $-106.96$ at the first coarse step, reference-side comparison), so this was
never the cause either.

\paragraph{The \texttt{grid\_scale} factor: real, but a reference-side deviation, not a paper
requirement.} The reference's trigger check, \texttt{c\_check}, does not compare
$\|R\operatorname{grad}f_h\|$ against $\eta\|\operatorname{grad}f_h\|$ directly -- it scales
the coarse side by a per-approach \texttt{grid\_scale} first
(\texttt{scaled\_norm\_coarse = grid\_scale * norm\_coarse}), and for Option 1 this factor is
$2^{n_{\text{pools}}}$ ($=2$ for one factor-2 step, \texttt{operators.py}). Eq.~(16) itself has
no such factor. \texttt{cc::ScaledCoarseCondition} (\S\ref{sec:condition-design}) implements it
as an explicit, labelled deviation, applied to \texttt{norm\_coarse} before both halves of
eq.~(16)'s comparison -- including the $\mu$ gate, which the same review (\S3.3) found the
port had been applying to the \emph{fine} norm instead of the coarse one as eq.~(16) states;
both variants give the same final energy on this problem, only the coarse-step count differs
(3 vs.\ 4).

\paragraph{$\varepsilon$ convention.} Eq.~(42) squares $\varepsilon$ under the root; the
reference's \texttt{objective.py} adds its \texttt{eps} argument unsquared. \S6.3.4 states the
paper's own figures were produced by the reference, so $\varepsilon_{\text{paper}} =
\sqrt{\varepsilon_{\text{ref}}}$ is the convention used from here on in both drivers:
$\varepsilon=10^{-2}$ fine (was $10^{-4}$), $\varepsilon\approx3.16\times10^{-2}$ coarse (was
$10^{-3}$). The previous values were the reference's own numbers passed unsquared into eq.~(42)'s
squared form, making the smoothing 100$\times$ sharper than anything the reference (and hence
the paper) ran; \texttt{REVIEW-continuous-cuts.md} \S3.3 measures the effect on the minimum
($-110.63$ vs.\ $-112.23$) and confirms it changes energies but not the qualitative behaviour
above.

\paragraph{Replicating the transfer operators outside deal.II.} \texttt{interpolate.h} and
\texttt{transport.h} were already free of deal.II's FE/mesh machinery -- \texttt{to\_fine\_mesh}
/\texttt{Tfine} are plain index-and-parity arithmetic, no \texttt{FE\_Q} or
\texttt{DoFHandler} involved, and their stencil coefficients were checked coefficient-for-
coefficient against \texttt{primitives.py}'s \texttt{prolong\_bilinear}/
\texttt{restrict\_adjoint\_bilinear} kernels (both golden-value-tested in
\texttt{test\_interpolate.cc}). The one remaining deal.II mesh dependency in the \texttt{cc}
module was \texttt{grid.h}'s \texttt{PixelGrid}, which built a throwaway
\texttt{Triangulation}+\texttt{DoFHandler}+\texttt{FE\_Q(1)} mesh solely to compute a
raster$\leftrightarrow$DoF permutation -- a leftover from an earlier plan (since abandoned) to
reuse deal.II's own mesh-transfer infrastructure for the grid operators. That permutation
turned out to always be the identity (deal.II's \texttt{FE\_Q(1)} support points on a
\texttt{subdivided\_hyper\_rectangle} enumerate in the same row-major order used everywhere
else in this module), and neither \texttt{get\_dofs()} nor \texttt{get\_triangulation()} were
called anywhere outside \texttt{grid.h} itself. \texttt{PixelGrid} was rewritten to plain
row-major indexing with no deal.II mesh includes at all (\texttt{to\_dof\_order}/
\texttt{to\_raster\_order} are now the identity); the \texttt{cc} module's only remaining
deal.II dependency is the linear-algebra containers (\texttt{Vector}, \texttt{SparseMatrix},
\texttt{DynamicSparsityPattern}), not FE/mesh infrastructure. The full \texttt{cc} test suite
and both drivers reproduce byte-identical results after the change.

\subsection{A pluggable coarse-correction condition}
\label{sec:condition-design}

The trigger condition (eq.~(16)) was hardcoded in \texttt{fas.h}'s \texttt{cycle()} as
\texttt{norm\_coarse >= options\_fas.kappa * norm\_level}. This is now factored out the same
way the framework already factors out the norm used to measure \texttt{norm\_level}/
\texttt{norm\_coarse} themselves (\texttt{LevelNorm}, \S\ref{sec:solveroptions}'s sibling
mechanism): a small interface,
\begin{quote}\ttfamily
class CoarseConditionBase \{ virtual bool trigger(double norm\_level, double norm\_coarse,
const FAS\_Options\&) const = 0; \};
\end{quote}
in \texttt{ropt/condition.h}, with \texttt{DefaultCoarseCondition} reproducing the original
formula exactly. \texttt{FullApproximationScheme} takes an optional
\texttt{MGLevelObject<shared\_ptr<CoarseConditionBase>>}, indexed the same way as
\texttt{cond\_norm\_mg}, with a null entry on any level falling back to
\texttt{DefaultCoarseCondition}; existing callers (\texttt{main\_coarse.cc}) that construct
\texttt{FullApproximationScheme} without this argument are unaffected. \texttt{cc/condition.h}
adds \texttt{ScaledCoarseCondition}, taking a \texttt{grid\_scale} constructor argument and
implementing \texttt{grid\_scale * norm\_coarse >= kappa * norm\_level} -- the reference's own
compensation from the previous section. \texttt{main\_cc\_coarse.cc} now installs
\texttt{ScaledCoarseCondition(2.0)} on the fine level, matching the reference's
\texttt{get\_grid\_scale("Option 1", 1)}; as predicted above, this changes nothing observable
(the trigger already fired on every iteration without it).

\section{Testing}
\label{sec:testing}

Four test files live in \texttt{test/cc}, one per header pair, named after what they test
(\texttt{test\_manifold.cc}, \texttt{test\_oracle.cc}, \texttt{test\_interpolate.cc},
\texttt{test\_oracle\_coarse.cc} -- no \texttt{cc\_} prefix, since the directory already
says that). Registered ctest names keep a \texttt{cc\_} prefix
(\texttt{cc\_manifold}, \texttt{cc\_oracle}, \texttt{cc\_interpolate},
\texttt{cc\_oracle\_coarse}) so a flat \texttt{ctest -N} listing still shows which module a
test belongs to, even though the file layout no longer needs to.

\texttt{test\_oracle\_coarse.cc} builds the full \texttt{CoarseOracleBase} machinery (a
$3\times3$ coarse / $5\times5$ fine grid, \texttt{GeometricBilinearTransport}) and checks, at a
random fine point, once for each of \texttt{BernoulliPointTransfer}'s two point restrictions
(injection: Table 10 Option 1; full-weighting: Option 4 -- they share $R,P$, differing only in
$r$, so this is also the first test exercising the full-weighting branch through the coherence
machinery rather than in isolation):
\begin{enumerate}[nosep]
  \item \textbf{Proposition 3.2 (first-order coherence)}: $q_k(\psi_k)=E_H(\psi_k)$ exactly,
        and $\operatorname{grad}q_k(\psi_k)$ equals the restricted fine gradient exactly.
        This is the strongest check available, since it is a property of the general
        framework rather than something derivable from the CC-specific formulas alone --
        it is what caught the gradient bug above.
  \item $\langle\operatorname{grad}q_k(z),v\rangle_z=\mathrm Dq_k(z)[v]$ at a point
        $z\neq\psi_k$ (via the oracle's own \texttt{inner}, not the plain dot product --
        \texttt{gradient()} returns the Riemannian gradient, so the plain dot product is
        the wrong thing to compare against \texttt{directional\_derivative}).
  \item A central-difference check of \texttt{directional\_derivative} against \texttt{value}.
  \item \texttt{residual(z) == norm(gradient(z))}.
\end{enumerate}

The same four-part pattern -- first-order coherence, the inner-product identity, a
central-difference check, and residual/gradient consistency -- was then written for GP's
own coarse oracle, in a new file \texttt{test/gpe/test\_coherence.cc}. It builds an actual
two-mesh-level GP problem (two independent \texttt{ModelBuilder}s, deal.II's
\texttt{LinearTransferMG} for the grid transfer, the mass metric throughout) and is, as far
as could be determined, the first test in this repository to exercise that transfer chain
end to end -- the existing GP tests check gradient formulas on a single mesh with a
synthetic tilt vector instead. It passes, including first-order coherence, confirming GP's
existing coarse-oracle code is correct on this path.

\subsection{Verification against the reference implementation}
\label{sec:verification}

\paragraph{The reference implementation.} The paper's authors also published a reference
implementation, \texttt{RMO-continuous-cuts} (python package \texttt{bernoulli\_multilevel},
built on \texttt{torch}), covering this same problem: the Fisher--Rao geometry, the grid
primitives, the Table 10 operator triples, and full single-/multi-level solvers. It is a
separate checkout, not part of this repository. Its modules, most to least relevant here:

\begin{center}
\begin{tabular}{@{}ll@{}}
\toprule
Module & Contents \\
\midrule
\texttt{manifold.py}       & Fisher--Rao metric: $G$, $G^{-1}$, the retraction (\texttt{exp}), lifting (\texttt{lifting}) \\
\texttt{primitives.py}     & factor-2 grid kernels: injection, bilinear interpolation, full-weighting \\
\texttt{objective.py}      & the smoothed-TV continuous-cuts objective \\
\texttt{operators.py}      & the five Table 10 \texttt{(r, R, P)} triples, their composition, and the registry \\
\texttt{optimizer.py}      & Riemannian gradient step + Armijo line search \\
\texttt{single\_level.py}  & single-level RGD (the paper's ``RGD'' baseline) \\
\texttt{multilevel.py}     & adaptive multilevel RGD and the coarse-correction trigger \\
\texttt{problem.py}        & builds a multilevel problem (and $\rho$) from an image and seed patches \\
\texttt{linear\_solvers.py}& the $(FF^\top)^{-1}$ solves Option 3's pseudoinverse needs \\
\bottomrule
\end{tabular}
\end{center}

Only the first four rows are used below, to derive golden values for \texttt{cc/manifold.h},
\texttt{cc/cc.h}, and \texttt{cc/interpolate.h}/\texttt{transport.h} respectively; the rest
describe the full reference solver (single- and multi-level), which \texttt{rmo::cc} does
not call into anywhere -- only its per-formula building blocks are used, as an independent
check.

A \emph{golden value} (or golden number) is a reference number computed once, by a trusted
and independent implementation, and hard-coded as the expected result in a test; the test
passes if the code under test reproduces it within a small tolerance. ``Golden'' just means
fixed and trusted, not that the test computes it itself.

Every golden value in \texttt{test\_manifold.cc}, \texttt{test\_oracle.cc}, and
\texttt{test\_interpolate.cc} was produced by calling the reference's own functions
directly on the fixed inputs below, so each check compares against the reference's actual
code, not a re-derivation of the same formulas. The script that does this is checked into
this repository as \texttt{test/cc/derive\_golden\_values.py}; running it reproduces every
constant used by the three test files:

\begin{lstlisting}[language=Python]
import sys
sys.path.insert(0, "src")

import torch
from bernoulli_multilevel.manifold import exp, lifting
from bernoulli_multilevel.objective import create_cc_objective
from bernoulli_multilevel.primitives import (
    r_injection, prolong_bilinear, restrict_adjoint_bilinear)
from bernoulli_multilevel.operators import BASE_TRIPLES

# -> test_manifold.cc
phi = torch.tensor([0.3, 0.6, 0.2, 0.8])
v   = torch.tensor([0.1, -0.2, 0.05, 0.3])
z = exp(phi, v)                    # -> retract(v, phi)
lifting(phi, z)                    # -> retract_inv(retract(v,phi), phi), should equal v

# -> test_oracle.cc
rho    = torch.tensor([[0.1, -0.2, 0.3], [0.0, 0.4, -0.1], [0.2, 0.1, -0.3]])
coarse = torch.tensor([[0.2, 0.5, 0.7], [0.3, 0.6, 0.1], [0.8, 0.4, 0.9]])
u0 = coarse.clone().requires_grad_(True)
E0 = create_cc_objective(rho, alpha=0.7, eps=1e-3)(u0)   # -> value(phi)
E0.backward()
u0.grad                            # -> gradient(phi), the ambient gradient

# -> test_interpolate.cc
fine_v = torch.arange(1., 26.).reshape(5, 5)
r_injection(fine_v)                          # -> to_coarse_mesh
prolong_bilinear(coarse, (5, 5))             # -> to_fine_mesh(coarse)
restrict_adjoint_bilinear(fine_v)            # -> Tfine(fine_v)

x_fine = torch.tensor([0.1 + 0.8*k/24.0 for k in range(25)]).reshape(5, 5)
v_coarse = torch.tensor([[1., 0., 2.], [0., 1., 0.], [2., 0., 1.]])
r1, R1, P1 = BASE_TRIPLES["Option 1"]
P1(x_fine, v_coarse, psi=coarse)             # -> vector_prolongation(...)

r4, R4, P4 = BASE_TRIPLES["Option 4"]        # confirms Options 1 & 4 share P, R:
torch.allclose(P1(x_fine, v_coarse, psi=coarse), P4(x_fine, v_coarse, psi=coarse))  # True
torch.allclose(R1(x_fine, fine_v, psi=coarse), R4(x_fine, fine_v, psi=coarse))      # True
\end{lstlisting}

Two things were caught this way, beyond confirming the formulas match: the Options 1/4
equivalence above (an earlier version of this plan had Option 4 wrong), and a units
mismatch in \texttt{objective.py}'s TV smoothing term, which adds \texttt{eps} (not
\texttt{eps}$^2$) under the square root, unlike eq.~(42)'s $\varepsilon^2$. \texttt{cc.h}
follows the paper ($\varepsilon^2$); the test compensates by passing $\sqrt{\texttt{eps}}$
as $\varepsilon$ so the two are compared on equal footing. This is a naming difference
between the two implementations, not a bug in either, but worth knowing if hyperparameters
are ever copied from one to the other.

\paragraph{The golden values as placeholders.} The literal numbers in the three test files
are not derived independently -- they are copies of \texttt{derive\_golden\_values.py}'s
output, kept in sync by hand. Each block is commented with which section of the script
produced it, so updating them is mechanical: if the reference implementation or a test's
fixed input changes, rerun the script and paste its new output over the corresponding
block. A stronger version of this -- a build step that runs the script and diffs its
output against the embedded literals automatically, so drift is caught rather than relied
on being noticed -- is possible but not implemented, since it would add a Python/\texttt{torch}
dependency to the C++ test suite for a check that, so far, only needs to run once per
change to either side.

\subsection{SolverOptions: an interface wrinkle}
\label{sec:solveroptions}

GP's oracles genuinely use \texttt{SolverOptions} (as tolerances for the inner CG solve);
CC's do not, since nothing in \texttt{rmo::cc} needs a linear solve. \texttt{FisherRaoOracle}
and \texttt{ContinuousCutsCoarseOracle} still take a \texttt{SolverOptions} parameter,
purely because the \texttt{TiltOracle}/\texttt{CoarseOracle} concepts
\texttt{FullApproximationScheme} relies on require a
\texttt{(Functional\&, SolverOptions)} (or \texttt{(CoarseOracleBase\&, SolverOptions)})
constructor uniformly, whether or not the metric needs one. This is not new: GP's own
\texttt{FrobeniusOracle} is in exactly the same position and already sets the precedent
(an ignored \texttt{SolverOptions} parameter with a default value). Left as-is for now --
two occurrences do not justify relaxing the concept, but if a third metric-with-no-solve
oracle shows up, it would be worth giving the concept a way to opt out of
\texttt{SolverOptions} instead of asking every such oracle to carry a dead parameter.

\section{Two pre-existing test failures, unrelated to this module}

\texttt{ctest} currently reports two failures outside \texttt{rmo::cc}: \texttt{arpack} and
\texttt{operator} (in \texttt{test/gpe} and \texttt{test/fe}). Neither file was touched
while building this module -- both were investigated because they showed up in the same
test run, not because anything here caused them.

\begin{description}
  \item[\texttt{operator}] crashes because \texttt{test\_operator.cc} calls
        \texttt{dof\_handler.distribute\_mg\_dofs()} on a plain
        \texttt{Triangulation} (default smoothing \texttt{none}), and deal.II's
        \texttt{distribute\_mg\_dofs()} requires the
        \texttt{limit\_level\_difference\_at\_vertices} smoothing flag to be set. This call
        turns out to be unnecessary: a standalone reproduction confirmed that
        \texttt{fe::LinearTransferMG} (the class the test exists to check) works correctly
        without it -- deal.II's two-DoFHandler, global-coarsening transfer only needs
        active degrees of freedom, not level-indexed ones. Removing the two
        \texttt{distribute\_mg\_dofs()} calls would very likely fix this, but that was not
        applied here since it was not asked for.
  \item[\texttt{arpack}] crashes because \texttt{test\_arpack.cc} deliberately sets
        \texttt{beta=0} (to reduce the Gross--Pitaevskii energy to a linear eigenvalue
        problem, the whole point of the test), and
        \texttt{GrossPitaevskiiFunctional}'s constructor calls
        \texttt{system.get\_operator\_A(beta)}, which adds the nonlinear term with weight
        \texttt{beta}. \texttt{LinearCombination::add\_component} asserts that a
        component's weight is non-zero, and this assertion is active in this Debug build.
        In a Release build the same zero-weight term would simply be a harmless no-op. So
        this is a Debug-only assertion that is stricter than the \texttt{beta=0} case this
        specific test deliberately exercises -- not a numerical bug, but a gap between a
        defensive check and a legitimate edge case.
\end{description}

Both were already present before this module's work began (neither file, nor
\texttt{lac.h}/\texttt{gpe.h} where the underlying checks live, appear in the changes made
here), and in both cases the mechanism, not just the symptom, has been identified above.

\section{Remaining work}

\begin{itemize}[nosep]
  \item Demonstrate an actual speedup from \texttt{main\_cc\_coarse.cc} over the
        single-level baseline, by tuning $\eta$ per problem or adding a
        \texttt{grid\_scale}-style compensation as discussed in
        \S\ref{sec:mlcc-design}. The wiring itself is done and correct.
  \item Once a CC-specific driver takes CLI options, add \texttt{FISHER\_RAO} to
        \texttt{option.h}'s parseable \texttt{MetricKind} list (\S 2).
\end{itemize}

\end{document}
